\section{Desenvolvimento da Solução}

Esta seção corresponde à etapa de \textit{Automação} do
ciclo de BPM apresentado na Base Teórica: o processo futuro
(TO-BE) especificado na seção de Análise e Redesenho do
Processo é aqui transformado em componentes executáveis do
ecossistema Microsoft 365, preservando, mesmo em ambiente
\textit{low-code}, as etapas de arquitetura, implementação e
validação próprias da Engenharia de Software. O
desenvolvimento foi realizado com base nos requisitos
consolidados na etapa anterior e utilizou componentes do
ecossistema Microsoft 365 já disponível na organização. A
solução resultante foi estruturada em cinco elementos
integrados: uma Central de Chamados no SharePoint, um
aplicativo Canvas no Power Apps para abertura das
solicitações, uma Microsoft List denominada
\textit{Chamados} como repositório central, dois fluxos
principais no Power Automate e uma área de acesso aos
indicadores gerenciais.

A arquitetura implementada difere parcialmente da
arquitetura inicialmente descrita na proposta. Na versão
desenvolvida, o SharePoint atua como portal de acesso,
enquanto o registro do chamado é realizado por um
aplicativo Power Apps incorporado à página. O Power Apps
grava diretamente o novo item na Microsoft List e, a partir
da criação desse item, o fluxo FL01 é disparado para gerar o
protocolo e realizar as notificações. Essa alteração
preservou o SharePoint como ponto único de entrada, porém
separou a camada de interface da camada de armazenamento,
permitindo maior controle sobre a experiência do
solicitante.

\subsection{Arquitetura da Solução Desenvolvida}

A Tabela~\ref{tab:arquitetura-desenvolvida} apresenta os
componentes efetivamente utilizados e suas responsabilidades
na versão desenvolvida.

\begin{table}[H]
\caption{Componentes da arquitetura desenvolvida}
\label{tab:arquitetura-desenvolvida}
\centering
\begin{tabular}{p{3.0cm} p{4.1cm} p{5.1cm}}
\hline
\textbf{Componente} & \textbf{Função principal} &
\textbf{Responsabilidade no fluxo} \\
\hline
SharePoint &
Portal corporativo da solução &
Centralizar o acesso para abertura de chamado, consulta de
registros e acesso ao painel gerencial. \\
\hline
Power Apps &
Interface de abertura &
Coletar os dados do solicitante, validar campos
obrigatórios, permitir anexos e gravar o novo chamado na
Microsoft List. \\
\hline
Microsoft List &
Repositório central &
Armazenar dados de abertura, classificação, responsável,
tratativa, status, datas e histórico de alterações. \\
\hline
Power Automate --- FL01 &
Automação da abertura &
Gerar protocolo, inicializar dados do chamado e notificar
Qualidade/SAC e solicitante. \\
\hline
Power Automate --- FL03 &
Automação da tratativa e encerramento &
Controlar notificações por mudança de status, transferência
de responsável, envio para análise, conclusão e
cancelamento. \\
\hline
Power BI &
Camada gerencial &
Disponibilizar a visualização dos indicadores e da base
histórica a partir do portal da Central de Chamados. \\
\hline
\end{tabular}
\raggedright
\par
{\fontsize{10}{12}\selectfont Fonte: Elaborado pelo autor}
\end{table}

O fluxo lógico da solução pode ser representado pela
sequência: solicitante acessa o SharePoint, realiza a
abertura no Power Apps, o registro é persistido na
Microsoft List, o FL01 gera protocolo e notificações, a
Qualidade/SAC classifica e direciona o chamado, a área
responsável realiza a tratativa e o FL03 controla as
mudanças subsequentes até o encerramento.

\subsection{Central de Chamados no SharePoint}

O SharePoint foi utilizado como camada de apresentação e
ponto central de acesso ao processo. A página desenvolvida
apresenta uma identificação institucional da Central de
Chamados Internos e disponibiliza acessos para abertura de
novo chamado, consulta dos chamados registrados e acesso ao
dashboard gerencial. Essa organização concentra em um único
ambiente as funções que, no processo anterior, estavam
distribuídas entre SAP Service, mensagens e comunicações
paralelas.

A utilização do SharePoint como portal também permite que o
processo seja incorporado ao ambiente corporativo já
utilizado pelos usuários, reduzindo a necessidade de
acesso a uma aplicação externa. O formulário de abertura
foi incorporado à página por meio de um aplicativo Power
Apps, enquanto a Microsoft List permanece responsável pelo
armazenamento e pelas etapas internas de classificação,
tratativa e encerramento.

\subsection{Aplicativo Power Apps para Abertura do Chamado}

A primeira etapa do formulário originalmente prevista
diretamente na Microsoft List foi substituída por um
aplicativo Canvas desenvolvido no Power Apps. A mudança foi
realizada com o objetivo de oferecer uma interface mais
simples ao solicitante e separar os campos de abertura dos
campos internos de classificação e tratamento.

O aplicativo foi integrado diretamente à lista
\textit{Chamados} e configurado para criação de novos
registros. Na versão analisada, o formulário de abertura
contém os campos Empresa, Tipo de chamado, Assunto do
chamado, Problema, Linha, Veículo e Anexos. Empresa, Tipo
de chamado, Assunto do chamado e Problema foram
configurados como obrigatórios. O campo Problema utiliza
entrada de texto em múltiplas linhas, possibilitando uma
descrição detalhada da ocorrência, enquanto o controle de
Anexos permite incluir evidências que apoiem a análise.

Os campos de escolha são carregados a partir das opções
cadastradas na própria Microsoft List, evitando a
manutenção de listas paralelas no aplicativo. O botão
\textit{Criar Chamado} utiliza a validação do formulário
antes do envio, enquanto o botão \textit{Limpar Formulário}
reinicializa os campos para uma nova entrada.

Após o salvamento bem-sucedido, o aplicativo altera seu
estado visual e apresenta ao usuário uma confirmação de que
o chamado foi aberto e está aguardando análise da
Qualidade/SAC. A mesma tela disponibiliza o botão
\textit{Novo Chamado}, que limpa o formulário, retorna ao
modo de criação e permite novo registro. Em caso de falha
no envio, uma mensagem de erro é apresentada ao usuário.

A interface foi ajustada para uso incorporado em página do
SharePoint, utilizando contêineres responsivos e dimensões
baseadas no espaço disponível. Essa adequação buscou
preservar a legibilidade do formulário quando executado
dentro do \textit{web part} do Power Apps.

\subsection{Estrutura da Microsoft List}

A Microsoft List denominada \textit{Chamados} constitui o
repositório central da solução. A estrutura separa os dados
informados pelo solicitante dos dados preenchidos durante o
processamento interno.

Na abertura são armazenados, entre outros, Empresa, Tipo de
chamado, Assunto, Problema, Veículo, Linha, anexos e o
usuário responsável pela criação do item. Nas etapas
posteriores são utilizados os campos Prioridade, Área
Responsável, Tipologia, Responsável pela tratativa, Causa
raiz, Ação corretiva, Data da Tratativa, Status do Chamado,
dados de transferência, Data de Fechamento e controle de
notificação de fechamento.

O registro centralizado permite que todas as etapas
utilizem o mesmo item como referência, evitando a criação
de bases paralelas e mantendo o histórico de alterações
associado ao protocolo. A estrutura também sustenta as
automações do Power Automate, que consultam alterações de
campos específicos para identificar quando uma nova ação
deve ser executada.

\subsection{Fluxo FL01 --- Abertura de Chamados Internos}

O fluxo FL01 é iniciado automaticamente quando um novo item
é criado na Microsoft List. Sua primeira ação calcula o
protocolo utilizando o ano corrente e o identificador
numérico do item. O padrão implementado é
\texttt{C-AAAA-NNNN}. Assim, um item de identificador 14
criado em 2026, por exemplo, recebe o protocolo
\texttt{C-2026-0014}.

Após a geração do protocolo, o próprio item é atualizado.
Na versão exportada do fluxo, o status inicial é definido
como \textit{Aberto} e a prioridade é inicialmente
registrada como \textit{Média}. Em seguida são executadas
duas notificações distintas.

A primeira notificação é destinada à equipe de
Qualidade/SAC e informa que um novo chamado aguarda
classificação. O e-mail apresenta protocolo, solicitante,
empresa, tipo, assunto, linha, veículo, status e descrição
do problema, além de um link direto para o item. O objetivo
é iniciar a etapa de classificação e direcionamento.

A segunda notificação é destinada ao próprio solicitante e
funciona como confirmação formal de abertura. O e-mail
informa o número do protocolo, resume os dados registrados,
indica que a solicitação aguarda análise da Qualidade/SAC e
fornece o link para acompanhamento do chamado. Com isso, a
confirmação de abertura deixa de depender de comunicação
manual.

\subsection{Classificação e Direcionamento pela
Qualidade/SAC}

Após a abertura, a Qualidade/SAC realiza a classificação
diretamente no registro da Microsoft List. Nessa etapa são
definidos os campos Prioridade, Área Responsável, Tipologia
e Responsável pela tratativa. O status do chamado também é
controlado pela Qualidade/SAC.

A decisão de manter esses campos fora do formulário do
solicitante reduz o risco de encaminhamento incorreto e
formaliza a Qualidade/SAC como instância de triagem. Quando
o chamado é validado e direcionado para tratamento, o
status é alterado para \textit{Em tratativa}, iniciando a
execução das regras do FL03 associadas à atribuição do
responsável.

A estrutura da lista também contempla o campo Motivo do
Cancelamento. Quando a Qualidade/SAC decide não prosseguir
com a solicitação e altera o status para
\textit{Cancelado}, o motivo registrado passa a compor a
rastreabilidade do encerramento sem tratativa.

\subsection{Fluxo FL03 --- Tratativa e Encerramento}

O FL03 é disparado quando um item da lista é modificado.
Inicialmente, o fluxo obtém quais propriedades foram
alteradas e recupera o estado atual do item. Em seguida,
uma estrutura de decisão do tipo \textit{Switch} avalia o
status do chamado e direciona o processamento para o ramo
correspondente.

No estado \textit{Em tratativa}, o fluxo verifica se houve
mudança do status para identificar uma nova atribuição.
Quando essa condição é atendida, o responsável pela
tratativa recebe uma notificação por e-mail contendo os
dados necessários e o link para acesso ao chamado.

Ainda no estado \textit{Em tratativa}, o fluxo verifica se
o campo de responsável foi alterado. Quando ocorre uma
transferência, o registro é atualizado com informações de
controle da transferência e o novo responsável é
notificado. Essa regra foi incorporada após a necessidade
identificada no levantamento de requisitos de redistribuir
chamados sem perda de rastreabilidade.

Quando Causa raiz e Ação corretiva são preenchidas e as
respectivas propriedades são identificadas como alteradas,
o FL03 atualiza automaticamente o status para
\textit{Em Análise} e envia uma notificação à
Qualidade/SAC. A área responsável, portanto, registra a
tratativa, mas não realiza o encerramento definitivo.

A análise final é responsabilidade da Qualidade/SAC. Após
avaliar a tratativa, a equipe pode alterar manualmente o
status para \textit{Concluído} ou, quando aplicável,
\textit{Cancelado}. No caso de conclusão, o FL03 notifica o
solicitante e os responsáveis envolvidos e registra a data
de fechamento.

Para impedir que a própria atualização realizada pelo
fluxo gere sucessivas notificações de encerramento, foi
utilizado um campo de controle denominado
``Notificação de fechamento enviada?''. O ramo de
conclusão somente envia o e-mail quando esse campo indica
que a notificação ainda não foi realizada; na sequência, o
campo é atualizado para \textit{Sim}. Esse mecanismo atua
como proteção contra reprocessamento e reduz o risco de
\textit{loop} provocado pelo gatilho de modificação da
lista.

No ramo \textit{Cancelado}, o solicitante recebe uma
notificação específica informando que o chamado foi
cancelado pela Qualidade/SAC. A comunicação apresenta os
dados do chamado e o motivo do cancelamento registrado no
item, mantendo o histórico disponível para consulta.

\subsection{Estados do Chamado e Responsabilidades}

A implementação consolidou cinco estados principais:
\textit{Aberto}, \textit{Em tratativa}, \textit{Em Análise},
\textit{Concluído} e \textit{Cancelado}. A
Tabela~\ref{tab:estados-chamado} apresenta a função de cada
estado.

\begin{table}[H]
\caption{Estados implementados no fluxo de chamados}
\label{tab:estados-chamado}
\centering
\begin{tabular}{p{2.6cm} p{4.2cm} p{5.5cm}}
\hline
\textbf{Status} & \textbf{Responsável predominante} &
\textbf{Significado no processo} \\
\hline
Aberto &
Qualidade/SAC &
Chamado registrado e aguardando classificação e
direcionamento. \\
\hline
Em tratativa &
Área responsável &
Chamado classificado, atribuído e disponível para registro
da causa raiz e da ação corretiva. \\
\hline
Em Análise &
Qualidade/SAC &
Tratativa preenchida pela área responsável e aguardando
validação antes do encerramento. \\
\hline
Concluído &
Qualidade/SAC &
Tratativa validada e chamado encerrado, com notificação aos
envolvidos. \\
\hline
Cancelado &
Qualidade/SAC &
Chamado encerrado sem continuidade da tratativa, com motivo
registrado e comunicação ao solicitante. \\
\hline
\end{tabular}
\raggedright
\par
{\fontsize{10}{12}\selectfont Fonte: Elaborado pelo autor}
\end{table}

\subsection{Rastreabilidade e Notificações}

A rastreabilidade foi implementada por meio da combinação
entre protocolo único, item persistido na Microsoft List,
histórico de modificações e notificações vinculadas às
mudanças de estado. Cada chamado permanece associado ao
mesmo registro durante todo o ciclo, permitindo recuperar
os dados de abertura, classificação, responsável, tratativa
e encerramento.

As notificações foram utilizadas como mecanismo de
transição entre os responsáveis pelo processo. Na abertura,
Qualidade/SAC e solicitante são comunicados. Na atribuição
ou transferência, o responsável pela tratativa é
notificado. Após o registro da causa raiz e da ação
corretiva, a Qualidade/SAC recebe a solicitação para
análise. Na conclusão, solicitante e responsáveis são
informados do encerramento; no cancelamento, o solicitante
recebe a justificativa correspondente.

Esse desenho substitui a comunicação paralela utilizada no
processo anterior por mensagens vinculadas a eventos
registrados na própria base, reduzindo a dependência de
grupos de mensagens para informar movimentações do chamado.

\subsection{Aderência da Implementação aos Requisitos}

A análise entre requisitos e componentes desenvolvidos
demonstrou que a versão atual atende aos elementos centrais
de centralização, protocolo, classificação, atribuição,
transferência, tratativa, validação e encerramento.
Entretanto, alguns requisitos identificados na pesquisa de
campo não puderam ser confirmados como implementados nos
artefatos analisados e devem ser tratados como itens de
evolução ou de validação antes da conclusão definitiva do
projeto.

\begin{table}[H]
\caption{Aderência dos principais requisitos à versão desenvolvida}
\label{tab:aderencia-requisitos}
\centering
\begin{tabular}{p{1.2cm} p{6.6cm} p{4.4cm}}
\hline
\textbf{ID} & \textbf{Requisito} &
\textbf{Situação observada} \\
\hline
RF01--RF14 &
Centralização, abertura, protocolo, notificações,
classificação, atribuição, transferência, tratativa,
análise e encerramento. &
Implementados nos componentes SharePoint, Power Apps,
Microsoft List, FL01 e FL03. \\
\hline
RF15 &
Controle automático de SLA por tipo e prioridade. &
Requisito identificado na pesquisa, porém o cálculo
automático não foi evidenciado nos fluxos exportados
analisados. \\
\hline
RF16 &
Comentários, atualizações e justificativas durante
tratativas prolongadas. &
Necessidade registrada pela Supervisora; não foi
identificado campo específico na estrutura final analisada. \\
\hline
RF17 &
Quantidade de protocolos relacionados. &
Requisito informado pela Analista; não foi identificado no
formulário final de abertura analisado. \\
\hline
RF18 &
Tipologias reutilizáveis. &
Implementado na etapa de classificação da Microsoft List. \\
\hline
RF19 &
Indicadores gerenciais. &
A Central de Chamados possui acesso ao dashboard; a
validação detalhada das métricas deve ser realizada na
etapa de avaliação dos resultados. \\
\hline
\end{tabular}
\raggedright
\par
{\fontsize{10}{12}\selectfont Fonte: Elaborado pelo autor}
\end{table}

Outro ponto de aderência parcial está na regra de
priorização. A pesquisa de campo indicou prioridade Alta
para chamados Massivos e Média para chamados Internos.
Entretanto, na versão exportada do FL01, todo novo chamado
é inicialmente atualizado com prioridade Média, ficando a
classificação definitiva sob responsabilidade da
Qualidade/SAC. Dessa forma, a regra de prioridade por tipo
não está automatizada no momento analisado.

Esses itens não invalidam o fluxo principal desenvolvido,
mas representam diferenças entre os requisitos levantados
e a versão implementada. O registro explícito dessas
diferenças é importante para a avaliação acadêmica da
solução, pois permite distinguir funcionalidades
efetivamente entregues de requisitos planejados ou
reservados para evolução.
